\chapter{Cryptosporidiosis}
\index{Cryptosporidiosis}

\section*{Definition and Overview}
Cryptosporidiosis is a diarrhoeal illness caused by the parasite \textit{Cryptosporidium}, a microscopic organism found in contaminated water, food, and soil. It is one of the most common causes of waterborne gastroenteritis worldwide, and it is highly contagious. In healthy people, infection usually causes watery diarrhoea that settles within one to two weeks. In people with weakened immune systems, children, and the elderly, the infection can be severe and prolonged. There is no specific cure, but most people recover with fluids and supportive care, and good hygiene prevents spread.

\section*{Epidemiology}
Cryptosporidiosis occurs worldwide and is a leading cause of waterborne disease outbreaks. In many countries, cases peak in summer and autumn, often linked to swimming pools, water parks, and contaminated drinking water. Young children are most commonly affected, and outbreaks occur in childcare centres. In developing countries, cryptosporidiosis is a major cause of diarrhoea and malnutrition in children. People with HIV and other immune deficiencies can develop chronic, severe infection that may be life-threatening.

\section*{Aetiology and Pathophysiology}
Cryptosporidiosis is caused by ingesting \textit{Cryptosporidium} oocysts, the hardy infective form of the parasite.

\begin{itemize}
    \item \textbf{Water:} swimming pools, lakes, rivers, and contaminated drinking water are common sources
    \item \textbf{Food:} contaminated food, unpasteurised milk, and unwashed produce
    \item \textbf{Person-to-person:} spread through poor hand hygiene, especially in childcare and households
    \item \textbf{Animal contact:} calves, lambs, and other animals can carry the parasite
    \item \textbf{Oocyst resistance:} the parasite survives in water and is resistant to chlorine
    \item \textbf{Infection:} the parasite invades the lining of the small intestine, causing inflammation and water loss
    \item \textbf{Immune clearance:} healthy immune systems clear the parasite within weeks; immune deficiency allows chronic infection
\end{itemize}

\section*{Clinical Features}
Symptoms begin one to twelve days after exposure and last one to two weeks.

\begin{itemize}
    \item Profuse, watery diarrhoea
    \item Abdominal cramps and pain
    \item Nausea and vomiting
    \item Low-grade fever
    \item Loss of appetite and weight loss
    \item Fatigue and dehydration
    \item Symptoms that can relapse just as they seem to settle
    \item In immunocompromised people: severe, prolonged, or chronic diarrhoea
\end{itemize}

\section*{Differential Diagnosis}
Other causes of acute watery diarrhoea are considered.

\begin{itemize}
    \item Viral gastroenteritis, including rotavirus and norovirus
    \item Bacterial gastroenteritis from salmonella, campylobacter, shigella, or E. coli
    \item Giardiasis and other parasitic infections
    \item Food poisoning
    \item Irritable bowel syndrome or other chronic bowel conditions
    \item Coeliac disease, when symptoms are prolonged
    \item Medication-related diarrhoea
\end{itemize}

\section*{When to Seek Medical Care}
See a doctor if diarrhoea is severe, prolonged, or accompanied by high fever, blood in the stool, or dehydration, or if the affected person is very young, elderly, or immunocompromised.

\textbf{Contact your local emergency services for:}
\begin{itemize}
    \item Signs of severe dehydration: dizziness, confusion, very little urine, or sunken eyes
    \item In a baby or child: not drinking, few wet nappies, lethargy, or dry mouth
    \item Severe abdominal pain or blood in the stool
    \item Vomiting that prevents keeping fluids down
\end{itemize}

\section*{Diagnosis and Investigations}
Diagnosis is confirmed by laboratory testing of stool.

\begin{itemize}
    \item \textbf{History:} symptoms, exposure to water, childcare, travel, and animal contact
    \item \textbf{Stool microscopy:} detection of oocysts under the microscope
    \item \textbf{Antigen testing:} a sensitive test for the parasite in stool
    \item \textbf{PCR testing:} molecular testing that can identify the species
    \item \textbf{Blood tests:} for dehydration, and for immune assessment in severe cases
    \item \textbf{Specialist referral:} for chronic infection in immunocompromised people
\end{itemize}

\section*{Management}
Treatment is supportive in most cases.

\begin{itemize}
    \item \textbf{Fluids:} oral rehydration solution or water with electrolytes to replace losses
    \item \textbf{Continue feeding:} babies should continue breast milk or formula
    \item \textbf{Anti-diarrhoeal medicines:} used with caution and usually not recommended in children
    \item \textbf{Nitazoxanide:} an antiparasitic medicine that may shorten illness in some people
    \item \textbf{Immune restoration:} for people with HIV, antiretroviral therapy is the key to clearing the infection
    \item \textbf{Hospital care:} intravenous fluids for severe dehydration
    \item \textbf{Staying home:} avoid childcare, school, and swimming while infectious
    \item \textbf{Notification:} cryptosporidiosis is a notifiable disease in many countries
\end{itemize}

\section*{Complications}
\begin{itemize}
    \item Severe dehydration, especially in infants and the elderly
    \item Electrolyte imbalance
    \item Weight loss and malnutrition
    \item Chronic or recurrent infection in immunocompromised people
    \item Spread to the biliary tract in severe immune deficiency
    \item Prolonged symptoms in some people
    \item Transmission to household and childcare contacts
\end{itemize}

\section*{Prognosis and Outcomes}
Healthy people usually recover fully within one to two weeks, although symptoms can persist longer and relapse is common. In immunocompromised people, particularly those with advanced HIV, the infection can be chronic, severe, and sometimes fatal, but restoration of the immune system with antiretroviral therapy usually clears the parasite. Children and the elderly recover well with careful attention to fluids. Good hygiene prevents reinfection and spread.

\section*{Prevention and Self-Care}
\begin{itemize}
    \item Wash hands thoroughly with soap and water after the toilet, nappy changes, and animal contact
    \item Do not swim for two weeks after diarrhoea stops; avoid swallowing pool water
    \item Drink treated or boiled water when travelling or where water safety is uncertain
    \item Wash fruit and vegetables and avoid unpasteurised milk
    \item People with immune deficiency should take extra care with water and hygiene
    \item Use alcohol-based sanitisers only as a supplement; soap and water are needed for cryptosporidium
    \item Stay home from childcare and school while symptomatic
\end{itemize}

\section*{Sources and Further Reading}
The following reputable sources provide further information for healthcare students and patients seeking reliable, current advice about cryptosporidiosis.

\begin{itemize}
    \item Healthdirect Australia. \textit{Cryptosporidiosis.} https://www.healthdirect.gov.au/
    \item Australian Government Department of Health and Aged Care. \textit{Cryptosporidiosis.} https://www.health.gov.au/
    \item Centers for Disease Control and Prevention. \textit{Cryptosporidium (crypto).} https://www.cdc.gov/
    \item World Health Organization. \textit{Cryptosporidiosis.} https://www.who.int/
    \item NSW Health. \textit{Cryptosporidiosis fact sheet.} https://www.health.nsw.gov.au/
    \item MSD Manual. \textit{Cryptosporidiosis: professional reference.} https://www.msdmanuals.com/
\end{itemize}
